% =====================================================================
%  CAPÍTULO 1 · INTRODUCCIÓN
%  Función: situar el problema, justificar por qué importa, fijar los
%  objetivos (numerados y medibles) y anunciar la estructura del documento.
%  Reglas clave (ver guía):
%   - Contexto (dónde estamos) ≠ Motivación (por qué hace falta resolverlo).
%   - La motivación NO adelanta resultados.
%   - Objetivos O1..On verificables: luego se cierran UNO A UNO en conclusiones.
%   - Define el nombre del producto y cada sigla en su primera aparición.
%   - Toda cifra o afirmación de mercado lleva cita.
% =====================================================================
\chapter{Introducción}
\label{cap:introduccion}

Este capítulo presenta el contexto del trabajo, la motivación que lo origina,
los objetivos que se persiguen y la organización del resto de la memoria.

\section{Contexto}
\label{sec:contexto}

Describe el ámbito en el que se sitúa el trabajo: qué sector, qué usuarios, qué
tecnología. Cada dato concreto necesita su fuente. Por ejemplo: según un informe
reciente, el 78\,\% de los estudiantes universitarios usa el móvil para
gestionar su actividad académica~\cite{informe_movil2025}.

% GUÍA: escribe la referencia pegada a la afirmación que sostiene, no al final
% del párrafo. Usa ~ antes de \cite para que no se separe de la palabra.

\section{Motivación}
\label{sec:motivacion}

Explica el problema concreto que se quiere resolver y por qué las soluciones
actuales no bastan. En este ejemplo, la reserva de salas de estudio obliga a
consultar varias fuentes y no informa de la ocupación real, un problema que
trabajos previos ya han señalado~\cite{garcia2024movil}. La aplicación que se
desarrolla en este \ac{TFG} se llama \emph{SalaLibre}.

\section{Objetivos}
\label{sec:objetivos}

El objetivo general de este trabajo es diseñar, implementar y validar una
aplicación móvil que permita consultar la ocupación de las salas de estudio en
tiempo real y reservarlas. Para alcanzarlo se plantean los siguientes objetivos
específicos:

\begin{enumerate}[label=\textbf{O\arabic*.}, ref=O\arabic*]
  \item \label{obj:estado} Analizar las soluciones existentes y sus carencias
        mediante una comparativa con criterios explícitos.
  \item \label{obj:app} Desarrollar una aplicación Android que permita reservar
        una sala en un máximo de dos pasos.
  \item \label{obj:rendimiento} Conseguir una latencia media inferior a
        \qty{100}{ms} con 200 usuarios concurrentes.
  \item \label{obj:validacion} Validar la usabilidad con usuarios reales y
        obtener una puntuación \ac{SUS} igual o superior a 70.
\end{enumerate}

% GUÍA: un objetivo es medible si al final puedes decir «cumplido / parcial /
% no cumplido» con una evidencia. «Mejorar el rendimiento» no lo es;
% «latencia media < 100 ms con 200 usuarios» sí.

\section{Alcance}
\label{sec:alcance}

Indica qué entra y qué queda fuera del trabajo. Por ejemplo, esta versión cubre
Android y la gestión de reservas, pero no incluye una versión iOS ni el pago de
servicios. Si durante el desarrollo el alcance cambia, actualiza esta sección y
los objetivos: el título y los objetivos no pueden prometer más de lo entregado.

\section{Estructura del documento}
\label{sec:estructura}

La memoria se organiza en siete capítulos. El Capítulo~\ref{cap:estado} revisa
los fundamentos y las soluciones existentes. El Capítulo~\ref{cap:metodologia}
describe la metodología, la planificación y la viabilidad del proyecto. El
Capítulo~\ref{cap:diseno} recoge los requisitos y el diseño de la solución, y el
Capítulo~\ref{cap:implementacion} su implementación. El
Capítulo~\ref{cap:pruebas} presenta las pruebas y los resultados obtenidos y,
por último, el Capítulo~\ref{cap:conclusiones} expone las conclusiones, las
limitaciones y el trabajo futuro. Los apéndices incluyen el enlace al
repositorio, el vídeo de demostración y el manual de instalación.

% GUÍA: esta lista debe coincidir EXACTAMENTE con el índice. Si cambias,
% añades o quitas un capítulo, actualiza este párrafo.
